\documentclass[11pt,a4paper]{article}
\usepackage[T1]{fontenc}\usepackage[utf8]{inputenc}\usepackage{mathptmx}
\usepackage[margin=2.3cm]{geometry}\usepackage{booktabs}\usepackage{parskip}
\title{\textbf{Résumé en français}\\\large Mesurer la résilience opérationnelle des écosystèmes FinTech en Afrique subsaharienne : le cadre AFRO-FINRES et l'application au Cameroun et à la CEMAC}
\date{}
\begin{document}\sloppy\maketitle
\section*{1. Problématique}
Les services financiers numériques (Mobile Money, paiements par API, agrégateurs) sont devenus une infrastructure critique en Afrique subsaharienne. Pourtant, aucun instrument comparatif ne mesure la capacité des écosystèmes FinTech nationaux à prévenir, absorber et surmonter une perturbation opérationnelle (cyberattaque, fraude, panne d'un prestataire, coupure d'électricité). Les cadres internationaux (Bâle 2021, NIST CSF 2.0, DORA, FSB 2023) sont conçus pour les économies avancées et pour les institutions individuelles.

\section*{2. Lacune de la littérature}
(i) La recherche sur la FinTech africaine porte surtout sur l'adoption et l'inclusion, pas sur la continuité des services. (ii) Les travaux sur la cyber-résilience financière sont conceptuels ou fondés sur des données confidentielles. (iii) Aucun indice ne combine cybersécurité, continuité, gouvernance et dépendances aux tiers. La CEMAC est particulièrement peu étudiée.

\section*{3. Questions de recherche et contribution}
\textbf{QR1} : comment conceptualiser et mesurer la résilience opérationnelle des écosystèmes FinTech avec des données publiques ? \textbf{QR2} : quelle est la situation de l'Afrique subsaharienne, de la CEMAC et du Cameroun ? \textbf{QR3} : la mesure est-elle robuste et associée aux usages de la finance numérique ?

La contribution est le cadre \textbf{AFRO-FINRES}, qui compte six dimensions : Prévention (P), Détection (D), Réponse (S), Reprise et continuité (C), Gouvernance (G) et Redondance des infrastructures tierces/partagées (T). Ces dimensions sont mesurées par 15 indicateurs publics : ITU GCI 2024, Banque mondiale WDI/WGI/Enterprise Surveys et PeeringDB.

\section*{4. Méthodologie}
L'étude suit une démarche de \emph{design science} (Peffers et al., 2007) et la méthodologie OCDE--JRC des indicateurs composites. Elle couvre 48 pays. Les indicateurs sont winsorisés (2,5--97,5\,\%), normalisés sur une échelle de 0 à 100 et agrégés avec des poids égaux. La robustesse est testée avec une agrégation géométrique, une agrégation modérée par la gouvernance, des poids ACP, 5\,000 simulations Monte Carlo et une analyse « leave-one-out ». La validation externe repose sur des corrélations de Spearman avec le Global Findex 2025, contrôlées par le PIB par habitant. Un module de stress-test illustratif (courbe de fonctionnalité de Bruneau et al., 2003) complète l'analyse avec quatre scénarios : fraude SIM-swap, panne d'API, panne de plateforme Mobile Money et attaque d'une infrastructure partagée.

\section*{5. Principaux résultats}
\begin{center}\begin{tabular}{lrrrrrrr}\toprule
 & P & D & S & C & G & T & Indice\\\midrule
Cameroun & 49,2 & 31,0 & 67,5 & 40,9 & 46,1 & 36,2 & \textbf{45,2} (25\textsuperscript{e}/48)\\
Moyenne CEMAC & 22,8 & 9,0 & 39,3 & 38,6 & 33,8 & 37,8 & 30,2\\
Moyenne ASS & 38,1 & 43,6 & 59,1 & 41,9 & 51,3 & 27,4 & 43,6\\\bottomrule\end{tabular}\end{center}
\begin{itemize}
\item Les scores vont de 3,6 (RCA) à 96,5 (Maurice). Le groupe de tête comprend Maurice, l'Afrique du Sud, le Botswana, le Ghana, le Rwanda, la Tanzanie, le Kenya et Cabo Verde.
\item La CEMAC a 13,4 points de retard sur la moyenne régionale. Son retard est surtout marqué en détection (9,0 contre 43,6).
\item Le principal point faible du Cameroun est la \textbf{détection} (pilier technique GCI : 6,18/20). Son point fort relatif est la réponse.
\item L'indice est cohérent (alpha de Cronbach = 0,86) et robuste (corrélations de rang de 0,97 à 0,99). L'intervalle Monte Carlo du rang du Cameroun va de 16 à 24.
\item L'indice est associé à la détention d'un compte ($\rho$ = 0,75) et d'un compte Mobile Money ($\rho$ = 0,45), y compris après contrôle du revenu.
\item Dans les stress-tests, porter la détection du Cameroun au niveau du quartile supérieur réduit de 19\,\% la perte simulée dans un scénario de fraude. Améliorer la gouvernance réduit la perte de 11,5\,\% en moyenne.
\end{itemize}

\section*{6. Implications et limites}
Les priorités pour le Cameroun et la CEMAC sont : renforcer les CIRT et le partage d'informations sur les menaces, améliorer la gouvernance et la supervision, et mettre en place un cadre régional de résilience opérationnelle (cartographie des services critiques, registre des prestataires, tests de scénarios). L'indice mesure une \emph{capacité} de résilience, pas la fraude ou les pannes réellement observées. Il est corrélé au niveau de développement, et l'analyse transversale ne permet pas de conclure à une causalité. Les paramètres des scénarios sont des hypothèses.

\section*{7. Revue cible recommandée}
\textbf{FinTech} (MDPI, ISSN 2674-1032) : indexée Scopus, ESCI et RePEc, en libre accès, avec des frais de publication (APC) de 1\,200 CHF (vérifiés en novembre 2026). Le format exigé est un résumé d'environ 200 mots et la structure Introduction, Matériels et méthodes, Résultats, Discussion, Conclusions, avec une déclaration de disponibilité des données obligatoire.
\end{document}
